\section{Attribute Selection for Fine-Grained Hierarchical Classification}

Attribute annotations provide useful semantic supervision for fine-grained recognition, but using all available attributes is not always optimal, since feature-selection studies emphasize the importance of removing irrelevant or redundant variables \cite{peng2005feature}, and recent work on attribute-based visual recognition also suggests that concise attribute subsets can preserve performance while improving interpretability \cite{yan2023learning}. In datasets such as CUB-200-2011, the attribute space is rich and detailed, covering color, shape, pattern, and body-part-related properties \cite{wah2011cub}. However, not every attribute is equally useful for distinguishing every bird species. Some attributes may be shared by many species, some may be redundant with other attributes, and some may be weakly related to the most difficult class distinctions. Therefore, when attributes are used as auxiliary supervision, attribute selection becomes an important design choice.

In general machine learning, feature selection aims to identify a subset of informative variables while removing irrelevant or redundant ones. Classical approaches often evaluate features according to their relevance to target labels and, in some cases, their redundancy with other features. For example, the minimum Redundancy Maximum Relevance (mRMR) criterion selects features that are highly relevant to the target while being minimally redundant with already selected features \cite{peng2005feature}. Relief-based methods estimate the quality of attributes by comparing nearby samples from the same and different classes, making them sensitive to local class distinctions \cite{kira1992practical, kononenko1994estimating}.In applied scorecard and risk modeling, Weight of Evidence and Information Value are commonly used during variable binning and screening to assess the usefulness of predictive variables \cite{siddiqi2012credit}.


Although these feature selection methods provide useful general principles, fine-grained hierarchical classification introduces additional considerations. First, the most useful attributes may not be globally discriminative across all classes. An attribute that separates visually unrelated birds may be less useful for distinguishing two sibling species within the same family. Second, the hierarchical structure of the labels means that discrimination can occur at multiple levels, such as order, family, and species. Third, in fine-grained bird classification, difficult errors often occur between visually and taxonomically related classes. Therefore, an effective attribute subset should support local discrimination among related classes, rather than only global separation across the entire dataset.


This idea is related to prior work on taxonomy-aware attribute querying. Mall et al. proposed a field-guide-inspired zero-shot learning framework in which a learner asks for the most useful attributes to describe a new class \cite{mall2021field}. Their method is motivated by real field guides, where a species can be described by a small number of informative attributes rather than requiring annotation of the complete attribute vector. In particular, their \emph{sibling-variance} strategy selects attributes that vary within a branch of the taxonomy, showing that taxonomy-aware attribute selection can be more informative than measuring attribute variation globally. This work is closely related to the motivation of this thesis, because both approaches emphasize attributes that distinguish related classes.


More broadly, attribute selection in this thesis is also related to hierarchical feature selection. Zhao et al. proposed hierarchical feature selection with recursive regularization, where different feature subsets are selected for different nodes in a class hierarchy by using parent-child and sibling relationships \cite{zhao2017hierarchical}. This line of work shows that hierarchical class structure can be useful for selecting features, because related nodes may share common features while sibling nodes may require discriminative features that separate them. However, these methods are general feature selection frameworks and are not specifically designed for CUB bird attributes or for attribute-guided visual representation learning.


This thesis considers attribute selection in the context of attribute-guided H-CAST. Since the proposed model uses attribute groups as auxiliary supervision and species-level feature guidance, the selected attributes influence which semantic cues are emphasized during training. If too many weak or redundant attributes are included, the attribute branch may learn less useful supervision and reduce the effectiveness of attribute-guided fusion. Conversely, if too few attributes are selected, the model may lose important cues needed to distinguish similar bird species. Attribute selection is therefore studied as a trade-off between compactness, discriminative power, and semantic coverage.

Several attribute usage strategies are considered in this thesis. The simplest strategy, referred to as \textbf{FullAttr-312}, uses all 312 CUB attributes. This setting preserves the full semantic information available in the dataset, but it may also include redundant or weakly discriminative attributes.

This thesis also includes an internal statistical attribute-selection baseline, referred to as \textbf{Sibling-IVAS-188}. This baseline was implemented as an exploratory comparison method rather than adopted from a published work. It uses train-only image-level binary attribute annotations and compares each taxonomy node with a corresponding negative set. For species-level nodes, the negative set consists of sibling species within the same family; for family-level nodes, it consists of sibling families within the same order; and for order-level nodes, it consists of the remaining orders. Attributes are ranked primarily by positive occurrence-rate difference between the target node and the negative set, with Information Value used as an additional ranking criterion. Therefore, Sibling-IVAS-188 should be understood as a taxonomy-aware statistical baseline, not as a global attribute-selection method.



To provide a stronger control, this thesis also includes \textbf{Random-199}, a random attribute subset with the same number of attributes as the main selected subset. This baseline is important because it helps determine whether any performance change is caused by the attribute selection criterion itself or merely by reducing the number of attributes. If a selected subset outperforms a random subset of the same size, this provides evidence that the selection strategy captures meaningful discriminative information.

The main attribute selection strategy studied in this thesis is \textbf{Hard-Sibling Discriminative Attribute Selection (HSDAS)}. HSDAS follows the high-level intuition of sibling-aware attribute selection: attributes should be chosen according to local taxonomic distinctions, not only according to global discriminative strength. However, the setting is different from field-guide-inspired zero-shot learning. Mall et al. use sibling-aware attribute querying to reduce the cost of expert annotation for novel classes \cite{mall2021field}. In contrast, this thesis assumes that class-level CUB attributes are already available during training and uses HSDAS to construct an attribute supervision mask for an H-CAST-based hierarchical classifier.

Conceptually, HSDAS reflects the nature of fine-grained bird classification. Two species from the same family may share global shape and overall color but differ in localized properties such as crown color, bill shape, wing pattern, or tail shape. Attributes that distinguish these hard sibling classes are more likely to support species-level recognition than attributes that separate visually unrelated birds. Therefore, HSDAS prioritizes attributes that distinguish a taxonomic node from its sibling nodes, especially hard siblings with similar attribute profiles.

In this thesis, HSDAS is not presented as a general-purpose feature selection algorithm for all datasets. Rather, it is designed for the specific setting where a taxonomy and class-level attribute annotations are available. Its role is to select a compact set of semantic attributes that can guide an H-CAST-based classifier toward discriminative bird traits. The selected attributes are then used to define active attribute groups for auxiliary attribute prediction and attribute-guided species-level fusion.

The study of attribute selection also supports the interpretability of the proposed approach. Since selected attributes correspond to human-readable bird traits, the resulting subset can be analyzed in terms of body parts and semantic groups. This makes it possible to inspect whether the model is guided by meaningful attributes such as bill shape, wing color, tail pattern, or breast markings. However, attribute selection does not imply that the model receives manually provided attribute labels at inference time. The attributes are used only during training as auxiliary supervision, while inference remains image-based.

Overall, attribute selection is a key component of the proposed attribute-guided hierarchical classification framework. It connects the semantic structure of CUB attributes with the taxonomic structure of bird labels. By comparing FullAttr-312, Sibling-IVAS-188, Random-199, and HSDAS-199, this thesis evaluates whether a compact and hard-sibling-aware attribute subset can better support fine-grained species recognition and hierarchical prediction than using all attributes, random subsets, or earlier internal attribute-selection baselines.
